\documentclass[../../../include/open-logic-section]{subfiles}

\begin{document}

\olfileid{sth}{replacement}{extrinsic}
\olsection[بهرني دليلونه]{د تعويض په اړه بهرني دليلونه}

لومړے د تعويض د توجيې يوه \emph{بهرنۍ} هڅه څېړو۔ بولوس ئې
داسې وړاندې کوي:
\begin{quote}
  [\ldots] د تعويض د اصولو د منلو دليل ډېر ساده دے: ګڼې غوښتل شوې
  پايلې لري او (ظاهراً) هېڅ ناغوښتلې پايله نۀ لري۔ د تکراري تصور په
  اړه د قضيو ترڅنګ، په دغو پايلو کښې د نامتناهي عددونو يوه د قناعت وړ،
  که څه هم نۀ بشپړه، تيوري او يوه ډېره مطلوبه پايله هم شامله ده چې په
  ښه‌بنسټيزو اړيکو باندې استقرايي تعريفونه توجيه کوي۔
  \citep[229]{Boolos1971}
\end{quote}
د بولوس د نظر لنډيز دا دے چې تعويض بايد د هغه د پايلو په مټ توجيه
کړو۔ هغه مشخصې پايلې همدا دي چې په تېرو څو څپرکو کښې مو وڅېړلې۔
تعويض موږ ته دا ثابته کړه چې د فون نويمن ترتيبي عددونه د ښه ترتيب د
ډول ډېر ښه استازي دي (دا زموږ «د نامتناهي عددونو د قناعت وړ، که څه
هم نۀ بشپړه، تيوري» ده)۔ تعويض مو د $V_\alpha$ ګانو تعريف، د رتبې
مفهوم جوړول، او د $\in$-استقرا ثبوت هم ممکن کړ (دا «د تکراري تصور په
اړه قضيې» دي)۔ بالاخره، تعويض د نامتناهيت هاخوا بازګښت د قضیې ثبوت
ممکن کړ (دا «په ښه‌بنسټيزو اړيکو باندې استقرايي تعريفونه» دي)۔

دا بېشکه غوښتل شوې پايلې دي۔ خو آيا همدغې پايلې د تعويض د
\emph{توجيې} دپاره بس دي؟ \emph{نۀ}۔ يا لږ تر لږه، نېغ په نېغه نۀ۔

يوه ساده ستونزه دا ده: که څه هم مو د تعويض ځينې مطلوبې پايلې
وښودلې، ډېرې ئې له نورو لارو هم ترلاسه کولے شو۔ دا خبره دومره مشهوره
نۀ ده څومره چې بايد وي؛ نو د حالت د روښانولو دپاره لږ تم کېږو۔

د سټونو يوه ساده تيوري شته چې د پړاوونو تيوري، يا په لنډه $\LT$
بلل کېږي۔\footnote{د $\LT$ لومړنۍ بڼې \citet{Montague1965} او
\citet{Scott1974} وړاندې کړې دي؛ \citet{Potter2004} ساده کړه او
د کتاب په کچه ئې وڅېړله؛ او \citet{ButtonLT1} په وروستيو کښې $\LT$
لا ساده کړې ده۔} د $\LT$ اصول يوازې د غړو له مخې برابري، بېلونه او دا
وينا دي چې هر سټ د کوم \emph{پړاو} فرعي سټ دے۔ «پړاو» په داسې کره ډول
تعريف شوے چې پړاوونه زموږ د $V_\alpha$ ګانو په شان چلند کوي۔ نو $\ZF$
د $\LT$ ثبوت ورکوي، خو $\LT$ له $\ZF$ څخه \emph{ډېره} کمزورې ده۔
په حقيقت کښې $\LT$ جوړې، د فرعي سټونو سټونه، نامتناهيت يا تعويض نۀ
راکوي۔ $\Zr$ دې د $\LT$ سره د نامتناهيت او د فرعي سټونو د سټونو د
اصلو د ورزياتولو پايله وي؛ له دې جوړې هم راځي، نو $\Zr$ لږ تر لږه د
$\Z$ هومره قوي ده۔ خو $\Zr$ په حقيقت کښې له $\Z$ څخه په کلکه قوي
ده، ځکه دا وينا ورزياتوي چې هر سټ يوه رتبه لري (له همدې امله ئې
$\Zr$ بولم)۔ په رښتيا $\Zr$ دا پايلې راکوي: د ترتيبي عددونو بشپړه د
قناعت وړ تيوري؛ هغه پايلې چې سلسله په ښه مرتب پړاوونو وېشي؛ د
$\in$-استقرا ثبوت؛ او د نامتناهيت هاخوا بازګښت يوه \emph{بڼه}۔

نو په لنډه توګه: که څه هم بولوس په دې نۀ پوهېده، هغه ټولې
مطلوبې پايلې چې يادوي ئې له تعويض \emph{پرته} هم ترلاسه کېدې؛ يوازې
$\Zr$ ئې د $\Z$ پر ځاے کارول پکار وو۔

(دې ټولو خبرو سره سره، ولې مو تاسو ته د $\ZF$ د ښوولو دوديزه لاره غوره کړه،
د $\LT$ او $\Zr$ د ښوولو پر ځاے؟ دوه دليلونه شته۔ لومړے، د
تاريخي لاملونو له امله له $\LT$ څخه پيلول ډېر نااشنا دي؛ غوښتل مو چې
تاسو د سټ تيورۍ نور دوديز بحثونه هم ولوستلای شئ۔ دويم، کله چې د
$\LT$ او $\Zr$ د درک دپاره چمتو شئ، د \citealt{Potter2004} او
\citealt{ButtonLT1} ليکنې پخپله لوستلے شئ۔)

البته، څرنګه چې $\Zr$ له $\ZF$ څخه په کلکه کمزورې ده، ځينې
داسې پايلې شته چې $\ZF$ ئې ثابتوي خو $\Zr$ ئې بې پرېکړې پرېږدي۔
نو څوک د دغو پايلو په يادولو د تعويض بهرنۍ توجيه هڅولے شي۔ خو کله
چې د $\Zr$ کارول زده کړئ، بيا د هغو خبرو ډېر مثالونه موندل ګران وي
چې (الف) تعويض ئې فيصله کوي او له هغه پرته نۀ فيصله کېږي، او (ب) په
بداهتي توګه رښتيا وي۔ (په دې اړه \citealt[\S13.2]{Potter2004}
وګورئ۔)

پايله دا ده: د تعويض د قناعت وړ بهرنۍ توجيې دپاره بايد داسې
پايله وموندل شي چې له تعويض \emph{پرته} نۀ شي ترلاسه کېدے۔ دا اسانه
کار نۀ دے۔

اوس يوه بله ستونزه وڅېړو چې د تعويض د هرې خالصې بهرنۍ توجيې
پر وړاندې درېږي۔ (ښايي دا له لومړۍ ستونزې هم ژوره وي۔) بولوس يوازې
دا نۀ وايي چې تعويض ډېرې مطلوبې پايلې لري؛ هغه دا هم وايي چې
«(ظاهراً) هېڅ ناغوښتلې» پايله نۀ لري۔ خو دلته د «ظاهراً» قيد بېخي
اساسي دے۔

راياد کړئ چې دلته څنګه راورسېدو: د ناپام جامعيت اصل له تناقض سره
مخ شو، او موږ په ځواب کښې د سټ تجمعي-تکراري تصور غوره کړ۔ له دې تصور
سره يوه کيسه تړلې ده چې هيله لرو د هغه سازګاري راته ډاډمنه کړي۔ خو
که تعويض له همدې کيسې نه توجيه نۀ کړو، تر اوسه د $\ZF$ د سازګارۍ
د باور دپاره دليل نۀ لرو۔ يا په کره توګه: د $\ZF$ د سازګارۍ د باور دپاره يوازې دا
(ښايي تصادفي) واقعيت راسره دے چې چا تر \emph{اوسه} تناقض نۀ دے موندلے۔ د بولوس خبره
په همدغه معنٰی دې ته راټيټېږي: «(ظاهراً) $\ZF$ سازګاره ده»۔ تر دې
د سازګارۍ ډېر ډاډ بايد وغواړو۔

دا ستونزه به د تعويض د هرې \emph{خالصې} بهرنۍ توجيې سره وي؛
يعنې له هغې توجيې سره چې يوازې د $\ZF$ په (معلومو) پايلو ولاړه وي۔
نو د تعويض \emph{دننۍ} توجيه غواړو: داسې توجيه چې وښيي د سټونو په اړه
زموږ کيسه موږ په يوه ډول تعويض ته «له مخکښې» ژمن کړي وو۔

\end{document}
